\section{Authoritative-state admission contract}
\label{sec:contract}

\subsection{Candidate, authorization, and authoritative state}


Correction-aware admission links an untrusted persisted candidate $c$, attributable authorization $a$, and authoritative pre-state $\Sigma_v$:
\begin{equation}
\operatorname{Admit}(c,a,\Sigma_v)=\Sigma_{v+1}.
\label{eq:admit}
\end{equation}
The candidate stays non-authoritative historical evidence. Proposal origin, value authority, and committed value remain distinct roles.

Let $r$ be a record with a nonempty authoritative field set $F_r$. Its committed state $\Sigma_v$ contains a total value map $V_v:F_r\rightarrow X$ and a total source map $S_v:F_r\rightarrow T$, where $T$ is the set of durable fact transitions. A machine candidate for field $f$ is
\begin{equation}
c=\langle r,f,e,v_{\mathrm{exp}},x_c,p,c_{\mathrm{id}}\rangle,
\label{eq:candidate}
\end{equation}
where $e$ is persisted evidence, $v_{\mathrm{exp}}$ is the expected current version, $x_c$ is the machine proposal, $p$ identifies its producer, and $c_{\mathrm{id}}$ is an identity over the modeled candidate content.


Evidence identity combines a content digest and canonical record/field locator:
\begin{equation}
EID(e)=\langle H(\mathrm{bytes}(e)),L(r,f,u,\kappa)\rangle.
\label{eq:evidence}
\end{equation}
Here $H$ is SHA-256, and $L$ contains the owner record, field, normalized storage-relative URI $u$, locator version, and optional crop $\kappa$. Alloy abstracts content and locators as primitive domains. Identity uniqueness in its scopes relies on collision resistance and canonicalization/store integrity in the implementation (Supplement S2).

An attributable authorization is
\begin{equation}
a=\langle c_{\mathrm{id}},x_a,q\rangle,
\label{eq:authorization}
\end{equation}
where $q$ is an authorized human principal and $x_a$ is the value explicitly permitted to enter the record. The dual-value semantics are
\begin{align}
\operatorname{proposal}(c)&=x_c, &
\operatorname{authorize}(a)&=x_a,\nonumber\\
\operatorname{commit}(\Sigma_{v+1},f)&=x_a, &
\operatorname{origin}(\Sigma_{v+1},f)&=c.
\label{eq:dual-value}
\end{align}
Accept has $x_c\eqc x_a$. Correction has $x_c\neqc x_a$; the candidate remains unchanged and the committed value derives its authority from $a$. Here $\operatorname{origin}$ identifies the retained proposal, not the authority for a corrected value.


\paragraph{Canonical value equality}
Write $x\eqc y$ when sorted-key, whitespace-minimal JSON serializations agree. This distinguishes $2$ from $2.0$ and $1$ from \code{true}, although Python equality does not. Change detection, value checks, tracing, the independent oracle, and projection use this relation. \Cref{sec:conformance} reports the cross-layer repair.

\subsection{History, observations, and outcomes}

An authority-relevant history is
\begin{equation}
h=\Sigma_0\xrightarrow{e_1}\Sigma_1\xrightarrow{e_2}\cdots\xrightarrow{e_n}\Sigma_n,
\end{equation}
where events may persist evidence or candidates, record review and authorization, attempt admission, commit or stutter, and reconstruct a reverse trace. The abstraction describes semantic information rather than requiring particular table or column names.

We group the observable distinctions into
\begin{equation}
\mathcal D=\{D_C,D_V,D_F,D_B,D_S\},
\label{eq:distinction-basis}
\end{equation}
where $D_C$ is candidate identity and context, $D_V$ is dual-value attribution, $D_F$ is freshness, $D_B$ is batch/successor completeness, and $D_S$ is total source attribution. For $I\subseteq\mathcal D$, $\pi_I(h)$ retains only the information classes in $I$. Candidate identity may be implemented by a certificate, stable handle, or another equivalence class, provided it does not collapse to value equality. Freshness may likewise use a version, state hash, epoch, or comparable discriminator.

To make ``retains'' precise, let $\operatorname{obs}_d(h)$ be the observation function for class $d$ and define
\begin{equation}
\pi_I(h)=\langle\operatorname{obs}_d(h)\rangle_{d\in I}.
\label{eq:observation-projection}
\end{equation}
The functions are logical projections, not claims that a database must store five physically separate objects. A certificate can co-locate several classes. In a projection, however, a content address is treated as an opaque equality-preserving handle: its preimage is not available as a side channel. This prevents a nominally omitted value or version from being recovered merely because the concrete certificate hash covers it. The complete observation functions are given in Supplement S2.

The full-history authority-outcome label is
\begin{equation}
O(h)=\langle status,successor,trace\rangle,
\label{eq:outcome}
\end{equation}
where $status$ is admissible, inadmissible, or failed/stuttering; $successor$ classifies the authoritative post-state or its absence; and $trace$ is complete, incomplete, ambiguous, or unavailable. This is the normative label assigned from the full history, not an additional input available to a mechanism restricted to $\pi_I$. A binary accept/reject label is insufficient because value-identical successors may differ in source completeness.

An observation set $I$ fails to distinguish a failure family when there exist a safe history $h_s$ and unsafe history $h_u$ such that
\begin{equation}
O(h_s)\neq O(h_u)
\quad\land\quad
\pi_I(h_s)=\pi_I(h_u).
\label{eq:indistinguishable}
\end{equation}
Any admission decision based only on that reduced observation cannot separate the pair.


The executable histories contain candidate and authorization registries, declared items, expected and observed versions, a connected commit sequence, endpoint maps, and transition sources. Projections and outcomes are derived from these relations. Schema fields determine total map domains; declared items determine required effects. Missing or conflicting sources remain failures in the domain. Supplement S2 gives the complete functions, finite tuple, and checker rules.

\subsection{Conditional failure-distinguishability characterization}

Table~\ref{tab:distinctions} summarizes the paired histories for this fixed observation basis.

\begin{table}[t]
\centering
\footnotesize
\caption{Failure-distinguishing information classes in the declared observation model.}
\label{tab:distinctions}
\begin{tabularx}{\textwidth}{@{}l>{\raggedright\arraybackslash}p{0.23\textwidth}>{\raggedright\arraybackslash}p{0.27\textwidth}>{\raggedright\arraybackslash}X@{}}
\toprule
Class & Safe/unsafe pair & Required distinction & Failure if omitted \\
\midrule
$D_C$ & exact candidate / equal-valued cross-context substitute & persisted candidate identity plus record, field, evidence, and producer context & candidate substitution \\
$D_V$ & preserved Correction / erased or false role attribution with the candidate unchanged & separate machine proposal $x_c$, human-authorized value $x_a$, committed value, and role relation & correction erasure or false machine attribution \\
$D_F$ & current authorization / same authorization after supersession & commit-time freshness discriminator tied to the observed pre-state & stale replay \\
$D_B$ & one atomic commit / two connected commits with the same final values & actual per-step effects and one atomic successor & fragmented successor \\
$D_S$ & total trace / equal-valued successor with missing or conflicting source & total per-field source relation and exact unchanged-source copy-forward & source ambiguity \\
\bottomrule
\end{tabularx}
\end{table}

\paragraph{Proposition 1 (conditional failure distinguishability)}
For each information class $d_i\in\mathcal D$, the declared failure model contains a safe history and a corresponding unsafe history whose authority outcomes differ but whose projections become equal when $d_i$ is omitted. Consequently, an admission mechanism that observes only $\mathcal D\setminus\{d_i\}$ cannot distinguish the paired failure family under this model.

\paragraph{Proof sketch (by construction)}
For $D_C$, hold the value-role, temporal, batch, and source observations constant while changing non-temporal candidate context. For $D_V$, retain the same candidate, proposal value 100, and committed value 101, but falsely attribute the corrected value to the machine proposal. For $D_F$, change only the admission-entry freshness comparison. For $D_B$, compare one two-field commit with two connected single-field commits that reach the same final values and sources; the latter exposes an intermediate partial state. Effect domains and successor counts follow from these sequences. For $D_S$, remove one source anchor while retaining candidates, values, freshness, and the commit sequence. Each construction changes its selected observation coordinate and the derived outcome, while leaving the other four coordinates equal.

Thus, for every $d_i\in\mathcal D$, the constructed histories satisfy
\begin{equation}
\pi_{\mathcal D\setminus\{d_i\}}(h_s^{(i)})
=
\pi_{\mathcal D\setminus\{d_i\}}(h_u^{(i)})
\quad\text{and}\quad
O(h_s^{(i)})\neq O(h_u^{(i)}).
\label{eq:irredundancy}
\end{equation}
Any decision function of the reduced projection receives the same input for both histories and cannot assign both distinct required outcomes. This establishes the claimed conditional irredundancy.

The executable checker validates each complete $\widehat h$ and derives observations and outcomes from its relations; witnesses store neither. It checks domain validity, unequal outcomes, equal four-class projections, and unequal full projections. Schema fields define the total value/source domains, whereas declared item fields define the required effect domain.

The original $D_C$ pair changes record context while retaining the candidate handle, establishing irredundancy of the combined class. An additional pair fixes the two-candidate registry and every context/value component, then selects either the authorized candidate or the equal-valued alternative. It isolates candidate binding within this observation model. A single-field copy-forward control checks the distinction between schema and effect domains. Supplement S2 identifies the executable constructions and raw histories.

\paragraph{Scope of the result}
Proposition 1 establishes class-wise irredundancy for the declared failure, history, observation, and outcome model, not universal minimality, schema uniqueness, or completeness over all admission failures. The checker validates the five class pairs, identity pair, and copy-forward control; it is not a general admission oracle. These constructions establish Eq.~\eqref{eq:irredundancy}; the separate Alloy checks test encoded assertions and conjunct-removal sensitivity in finite scopes.

\subsection{Batch admission relation}

A batch $B$ is a nonempty set of items sharing one target record, principal, expected pre-version, and successor. Under Full, item fields are unique. Let
\begin{equation}
\Delta(V_v,V_{v+1})=\{f\in F_r\mid V_v(f)\neqc V_{v+1}(f)\}.
\end{equation}
For a post-initial concrete admission,
\begin{equation}
\Delta(V_v,V_{v+1})=\{i.f\mid i\in B\}=\{t.f\mid t\in T_{v+1}\},
\label{eq:changed-fields}
\end{equation}
with one transition per item. At the first certificate-era version, the batch covers all of $F_r$, establishing total value and source domains. Deletion is outside the model.

For every $i\in B$, admission requires
\begin{align}
i.r &= r, & i.f&\in F_r,\nonumber\\
EID(i.e) &= EID(i.c.e), & i.v_{\mathrm{exp}}&=v,\nonumber\\
i.a.c_{\mathrm{id}}&=i.c.c_{\mathrm{id}}, & i.a.q&=q,\nonumber\\
i.x &\eqc i.a.x_a, & t_i.x&\eqc V_{v+1}(i.f)\eqc i.x.
\label{eq:item-bindings}
\end{align}
Together these instantiate the admission preconditions
\begin{equation}
\begin{aligned}
&\operatorname{Bind}(a,c)\land\operatorname{Fresh}(c,\Sigma_v)\land\operatorname{Authorized}(a)\\
&\qquad\land\operatorname{Complete}(\Sigma_{v+1})\land\operatorname{Attributed}(\Sigma_{v+1}).
\end{aligned}
\end{equation}

The successor is atomic: all admission-time decision, authorization, transition, version, source-map, and audit effects commit together, or authoritative state stutters. This does not require rollback of evidence, candidates, or review intent persisted before the admission attempt. For a changed field, $S_{v+1}(f)=t_i$. For an unchanged field, value and exact source are copied forward:
\begin{equation}
f\notin\Delta(V_v,V_{v+1})\Rightarrow
V_{v+1}(f)\eqc V_v(f)\land S_{v+1}(f)=S_v(f).
\label{eq:copy-forward}
\end{equation}
Alloy permits same-value admission for correspondence with the historical singleton model; concrete post-initial admissions require a value change. For submitted items $U$, the admitted batch is $B=\{i\in U\mid i.x\neqc V_v(i.f)\}$, with $B\neq\varnothing$. A mixed submission admits exactly $B$ and copies unchanged items' exact prior sources into the complete successor, creating no transition or authorization binding for $U\setminus B$. Thus Eq.~\eqref{eq:changed-fields} concerns admitted effects, not all submitted reviews. Recording renewed review or re-attribution without a value change requires a separate review-event path outside this contract; persisted candidate identity allows that path to distinguish equal-valued proposals.

\subsection{Operational properties and trust boundary}

\begin{table}[t]
\centering
\caption{Admission properties and their operational interpretation.}
\label{tab:properties}
\begin{tabularx}{\textwidth}{@{}l>{\raggedright\arraybackslash}p{0.29\textwidth}>{\raggedright\arraybackslash}X@{}}
\toprule
Property & Requirement & Operational interpretation \\
\midrule
P0 & Direct-write exclusion & Machine/candidate capability cannot change authoritative values, versions, sources, transitions, or authorizations. \\
P1 & Candidate non-substitutability & Authority refers to the same canonical persisted certificate, not merely an equal value or analogous candidate. \\
P2 & Context-bound admission & Record, field, persisted evidence content and locator, and expected version agree across the item and certificate. \\
P3 & Authorization integrity & The principal is allowed at the in-transaction policy check before the compare-and-swap, and the authorization names the exact certificate and explicit committed value; in-flight revocation semantics are stated in \cref{sec:authorization-timing}. \\
P4 & Freshness & The expected version equals the transactionally observed current version; stale or competing decisions fail closed. \\
P5 & Successor integrity & One atomic successor has an item--transition bijection, exact changed-field effects, and complete changed/unchanged framing. \\
P6 & Trace soundness & Each committed field resolves through its exact source transition to consistent authorization, certificate, evidence, producer, and version relations. \\
\bottomrule
\end{tabularx}
\end{table}


Model output and candidate-only callers are untrusted. The trusted base comprises the admission service, adapter, exercised transaction semantics, canonicalization and hashing, and principal policy. Host authentication, session integrity, and faithful review presentation are external assumptions. The confirmation API constructs its binding from caller-supplied candidate and value; the contract preserves that binding, not factual truth or protection from a compromised host.

\paragraph{Authorization timing and revocation}
\label{sec:authorization-timing}
The in-transaction principal check precedes compare-and-swap and authority writes. Revocation visible before the check causes zero-write rejection; a later overlapping revocation does not cancel an in-flight commit. This is commit-entry revalidation. Supplement S2 states the policy-lock and stronger-revocation boundaries.
